\section{Integración híbrida y contratos de frontera}

Los contratos de frontera rigen las interfaces entre módulos y constituyen una especificación distinta de la organización física del repositorio. El repositorio indica dónde se conserva cada componente; los contratos indican cómo puede interactuar.

\begin{table}[H]
\centering
\caption{Fronteras funcionales principales.}
\begin{tabular}{p{0.27\textwidth}p{0.60\textwidth}}
\toprule
Frontera & Responsabilidad \\
\midrule
Referencia $\rightarrow$ FOC & referencias y límites contractuales \\
FOC $\rightarrow$ M3 & $v_d$, $v_q$ y señales requeridas \\
M3 $\rightarrow$ M2 & señales funcionales de actuación \\
M2 $\rightarrow$ M4 & acción equivalente aplicada a la planta \\
M4 $\rightarrow$ M1/M6 & corrientes, posición y velocidad \\
M7 $\rightarrow$ M6 & información predictiva/adaptativa acotada; no sustituye al FOC \\
\bottomrule
\end{tabular}
\end{table}

La arquitectura híbrida deberá preservar tres principios: el FOC mantiene el control directo; la inteligencia actúa mediante contratos explícitos; y toda función inteligente puede desarrollarse y verificarse primero en Python antes de su implementación en Simulink/Simscape.
